% -*- coding: utf-8 -*-
% !TEX program = xelatex
\documentclass[plain,noanswer,medmath]{../../template/jnuexam}
\usepackage{../../template/kysx}

\begin{document}


\examtitle{2010}{3}


\loadproblems{1}{2010P1}%载入试卷一
\loadproblems{2}{2010P2}%载入试卷二


\makepart{选择题}{1～8小题，每小题4分，共32分}


\begin{problem}
若$\lim_{x \to 0} \left[\frac{1}{x} - \left(\frac{1}{x} - a \right) \e^x \right] = 1$，则$a$等于\pickout{}
\begin{abcd}
\item $0$．
\item $1$．
\item $2$．
\item $3$．
\end{abcd}
\end{problem}


\useproblem{2}{1}{2}%【同试卷二第一(2)题】


\begin{problem}
设函数$f(x),g(x)$具有二阶导数，且$g''(x) < 0$，若$g(x_0) = a$是$g(x)$的极值，
则$f(g(x))$在$x_0$取极大值的一个充分条件是\pickout{}
\begin{abcd}
\item $f'(a) < 0$．
\item $f'(a) > 0$．
\item $f''(a) < 0$．
\item $f''(a) > 0$．
\end{abcd}
\end{problem}


\begin{problem}
设$f(x) = \ln^{10} x$，$g(x) = x$，$h(x) = \e^{\frac{x}{10}}$，则当$x$充分大时有\pickout{}
\begin{abcd}
\item $g(x) < h(x) < f(x)$．
\item $h(x) < g(x) < f(x)$．
\item $f(x) < g(x) < h(x)$．
\item $g(x) < f(x) < h(x)$．
\end{abcd}
\end{problem}


\useproblem{2}{1}{7}%【同试卷二第一(7)题】


\useproblem{1}{1}{6}%【同试卷一第一(6)题】


\useproblem{1}{1}{7}%【同试卷一第一(7)题】


\useproblem{1}{1}{8}%【同试卷一第一(8)题】


\makepart{填空题}{9～14小题，每小题4分，共24分}


\begin{problem}
设可导函数$y = y(x)$由方程$\int_0^{x + y} \e^{- t^2} \dt = \int_0^x x\sin t^2 \dt$确定，
则$\left. \frac{\dy}{\dx} \right|_{x = 0} =$ \fillin{}．
\end{problem}


\begin{problem}
设位于曲线$y = \frac{1}{\sqrt{x(1 + \ln^2 x)}}$（$\e \le x < +\infty$）下方，
$x$轴上方的无界区域为$G$，则$G$绕$x$轴旋转一周所得空间区域的体积是 \fillin{}．
\end{problem}


\begin{problem}
设某商品的收益函数为$R(p)$，收益弹性为$1 + p^3$，其中$p$为价格，且$R(1) = 1$，则$R(p) =$ \fillin{}．
\end{problem}


\begin{problem}
若曲线$y = x^3 + ax^2 + bx + 1$有拐点$(- 1,0)$，则$b =$ \fillin{}．
\end{problem}


\useproblem{2}{2}{14}%【同试卷二第二(14)题】


\begin{problem}
设$X_1,X_2, \ldots ,X_n$是来自总体$N(\mu ,\sigma^2)$$(\sigma > 0)$的简单随机样本，
统计量$T = \frac{1}{n}\sum_{i = 1}^n X_i^2$，则$E(T) =$ \fillin{}．
\end{problem}


\makepart{解答题}{15～23小题，共94分}


\begin{problem}[points=10]
求极限$\lim_{x \to +\infty} (x^{\frac{1}{x}} - 1)^{\frac{1}{\ln x}}$．
\end{problem}


\begin{problem}[points=10]
计算二重积分$\iint_D (x + y)^3 \dx\dy$，
其中$D$由曲线$x = \sqrt{1 + y^2}$与直线$x + \sqrt{2} y = 0$及$x - \sqrt{2} y = 0$围成．
\end{problem}


\begin{problem}[points=10]
求函数$u = xy + 2yz$在约束条件$x^2 + y^2 + z^2 = 10$下的最大值和最小值．
\end{problem}


\useproblem[points=10]{1}{3}{17}%【同试卷一第三(17)题】


\begin{problem}[points=10]
设函数$f(x)$在$\left[0,3 \right]$上连续，在$\left(0,3 \right)$内存在二阶导数，且
$$2f(0) = \int_0^2 f(x)\dx = f(2) + f(3).$$
\begin{enumerate*}
\item 证明存在$\eta \in(0,2)$，使$f(\eta) = f(0)$；
\item 证明存在$\xi \in(0,3)$，使$f''(\xi) = 0$．
\end{enumerate*}
\end{problem}


\useproblem[points=11]{1}{3}{20}%【同试卷一第三(20)题】


\useproblem[points=11]{2}{3}{23}%【同试卷二第三(23)题】


\useproblem[points=11]{1}{3}{22}%【同试卷一第三(22)题】


\begin{problem}[points=11]
箱中装有$6$个球，其中红、白、黑球的个数分别为$1,2,3$个，
现从箱中随机取出$2$个球，记$X$为取出的红球个数，$Y$为取出的白球个数．\par
\begin{enumerate*}
\item 求随机变量$(X,Y)$的概率分布； 
\item 求$\Cov (X,Y)$．
\end{enumerate*}
\end{problem}


\end{document}
